\documentclass[11pt]{article}
\usepackage[margin=1.1in]{geometry}
\usepackage{amsmath,amssymb,amsthm}
\usepackage{booktabs}
\usepackage{array}
\usepackage{tabularx}
\usepackage[protrusion=true,expansion=false]{microtype}
\usepackage[table]{xcolor}
\usepackage{enumitem}
\usepackage{titlesec}
\usepackage{float}
\usepackage{hyperref}

\hypersetup{colorlinks=true, linkcolor=blue!60!black, urlcolor=blue!60!black, citecolor=blue!60!black}
\usepackage[skip=0.6\baselineskip, indent=0pt]{parskip}

\titleformat{\section}{\large\bfseries}{\thesection.}{0.7em}{}
\titleformat{\subsection}{\normalsize\bfseries}{\thesubsection.}{0.6em}{}

\newcommand{\vmax}{V_{\max}}
\newcommand{\Tc}{T_{\mathrm{close}}}
\newcommand{\bcorr}{b_{\mathrm{corr}}}
\newcommand{\tx}[1]{\texttt{#1}}
\newcolumntype{L}{>{\raggedright\arraybackslash}X}

\renewcommand{\arraystretch}{1.3}
\makeatletter
\setlength{\@fptop}{0pt}
\setlength{\@fpsep}{1.5\baselineskip}
\setlength{\@fpbot}{0pt plus 1fil}
\makeatother
\rowcolors{1}{}{black!5}

\title{\textbf{The grants L2}\\a Lichen L2 without global order\\[2pt]\large early draft}
\author{waxwing}
\date{Version 0.1, November 2026\\[2pt]\small a companion to \emph{Lichen}, version 0.5~\cite{lichen}}

\begin{document}
\maketitle

\begin{abstract}
Lichen~\cite{lichen} enforces an L2's withdrawals through bilateral
contracts between a hub and each user, settled on Bitcoin by BABE with
the roles swapped: the user claims, and only the hub proves, if it can,
that she is owed less. This draft describes a second L2 for those
contracts, with no sequencer, no global ordering and no anchoring. Notes
are transparent and signed by the hub; a payment is made atomic by a
\emph{grant}, an adaptor signature that the payer's own signature
releases, as in the ecash design of~\cite{ecash} without the blinding.
An exit posts a compact claim of the user's notes, one-time signed in
full. The hub answers a claim it can show is too high by posting
$\bcorr$, what she is owed, and proves it from spends she signed. A
spend counts against her only if the hub reveals its unlock to her on
chain, which the hub can do only by completing, in its own transaction,
a pre-signature she checked before paying. We give the algorithm, the
transaction graph, the statement the hub proves, and the costs measured
on regtest: an exit's \tx{request} is 11.7~kvB, the hub's correction
0.65~kvB, and the statement 1.3M RISC Zero cycles, 2 segments.
\end{abstract}

\section{Setting}

A Lichen contract (\cite{lichen}, \S3) locks the hub's liquidity and the
user's deposit in one Bitcoin output, with a cap $\vmax$ in a unit $u$
($b$ units claimable, $b < 2^{k}$), reserves $r_A$ and $r_H$, and a
close time $\Tc$. It ends in one of three ways:
\begin{itemize}[nosep]
\item \textbf{Fold}: both sign; the user is paid what she withdraws and
  her reserve.
\item \textbf{Default}: after $\Tc$, the user is paid her reserve and the
  hub the rest. She must withdraw or exit before.
\item \textbf{Exit}: the user posts \tx{request} with a claim of $b$
  units; the hub either lets it pay or proves her owed less.
\end{itemize}

The anchored L2 (the reference L2 of~\cite{lichen}) judges a claim
against a global, anchored history. The grants L2 drops that: the hub's
book is the only ledger, and nothing is ordered beyond what each spend
names. What replaces the history is \emph{evidence the user signed}: a
claim is reduced only by spends she authorized whose completion she was
given.

\textbf{Contracts and credit.} A user's position on the L2 is defined by
her contract. Each contract has an L2 key of its own, fresh at the
opening; the notes owned by that key are what its exit can claim. Notes
held by any other key (a wallet without a contract, a closed contract's
key) are the hub's IOUs with no cryptographic backing: credit. A wallet
moves its notes to its open contract's key with an ordinary payment to
itself.

\section{Notes and payments}

\subsection{Notes}

A \emph{note} is $(\mathit{id}, \mathit{owner}, v)$: a 32-byte id, an
owner's x-only key, a value in sats. It is \emph{settled} when its holder
has the hub's BIP340 signature, under the issuer key $I$, on
\[
  \mathrm{digest} = H_{\text{note}}(\mathit{id}_{[0..8]} \,\|\,
  \mathit{owner} \,\|\, v),
\]
the first 8 bytes of the id standing for it (an exit names a note by
them, \S\ref{sec:claim}). The book keys notes by the full id. Two notes
of one owner with the same short id and value would share a signature;
only their owner can grind for that ($2^{32}$ spends) and it only harms
her.

\subsection{Spends and grants}\label{sec:grant}

A \emph{spend} names input notes and asks for outputs $(o_j, v_j)$ of the
same total; output $j$'s id is $H(\mathit{sid} \,\|\, j)$, where
$\mathit{sid}$ hashes the inputs and outputs. No two outputs go to one
owner (\S\ref{sec:volume}). A payment is one round trip:

\begin{enumerate}[nosep]
\item \textbf{Grant} (\tx{POST /grant}). The payer sends the unsigned
  spend. The hub, if its inputs are live, answers with
  \begin{itemize}[nosep]
  \item $T = tG$, signed with $\mathit{sid}$, where
    $t = H(k_{\text{iss}}, \mathit{sid})$;
  \item an adaptor pre-signature under $T$ on each output note's digest;
  \item a \emph{reveal} for each input owner with a contract open: an
    adaptor pre-signature $\sigma' = (R', s')$ under $T$, by the hub's
    key $P_H$ in that contract, on the contract's \tx{correct}
    sighash~$m_c$ (\S\ref{sec:reveal}).
  \end{itemize}
\item \textbf{Check and sign.} The payer verifies $T$'s signature and
  every pre-signature, the reveal against her own contract's $m_c$ in
  particular. She signs $H(\mathit{sid} \,\|\, C)$, where
  \[
    C = H_{\text{commit}}(\mathit{sid} \,\|\, T \,\|\, \text{output
    pre-signatures} \,\|\, \text{reveals})
  \]
  is the grant's commitment, with each input's owner key.
\item \textbf{Record} (\tx{POST /spend}). The hub checks the signatures,
  marks the inputs spent, enters the outputs in the book and in their
  owners' mailboxes, and answers with $t$.
\end{enumerate}

With $t$, every output's pre-signature completes to the hub's ordinary
signature: the outputs are settled. Before $t$ they are unusable. The
payer keeps the spend \emph{in flight} until she has $t$. A hub that
says it never granted a spend is believed for spending its inputs again,
but not for claiming them: the spend remains a \emph{fallback} of any
claim of its inputs until a spend she holds $t$ for consumes them.

\textbf{What atomicity means here.} Once the payer signs, the hub may
record the spend or not. If it records it and gives $t$, both sides are
paid. If it records it and withholds $t$, it holds her inputs and pays
no one; the payee has nothing to deliver goods against. The exit
(\S\ref{sec:exit}) settles that case: the hub can use the spend against
her claim only by revealing $t$ to her, after which she completes the
outputs and hands the payee his note. Otherwise the spend doesn't count,
and her claim is owed its inputs in full.

\subsection{The wallet's policy}

A payment spends all the wallet's notes (the payment out, one change
note back). Received notes are merged by a payment to itself once there
are more than 2; more than 4 unmerged means the hub refuses merges, and
the wallet should exit while every note still fits in a claim. The
wallet's address is its open contract's key.

\section{The contract}

\subsection{Opening}

The opening follows~\cite{lichen}: the user's deposit to a 2-of-2 with
a refund after a timeout; the setup (the hub's statement fixed, its
journal below; the user garbles the BABE instances, the hub keeps some
by cut-and-choose); the pre-signed graph (\S\ref{sec:graph}); the
funding, which spends the deposit with the user's adaptor signature
under the hub's point. The user's deposit is \emph{pre-issued} to her as
a note under the same point: the funding on chain reveals $t$ and the
note settles, so the deposit and its L2 balance swap atomically.

The hub refuses a contract whose owner key it has seen before (a
contract of it, or any note), so a contract's key has no spends from
before its opening (\S\ref{sec:reveal}).

\textbf{The journal}, the statement's only public input, fixed at the
opening:
\[
  c \,\|\, D \,\|\, H(O_D \,\|\, \mathit{margin}),\qquad
  D = H\big(H(I \,\|\, \mathit{owner}) \,\|\, 0 \,\|\, W \,\|\, u\big),
\]
with $c$ the contract number, $O_D$ the deposit coin, $W$ the work the
statement asks on top of \tx{request}'s block (ten blocks' at the
opening's difficulty), $u$ the unit. $D$ binds the issuer and the
owner: without the owner, a hub could prove against another key and
count any claim malformed.

\subsection{Withdrawal}

The user hands the hub a claim (\S\ref{sec:claim}) signed with her
contract key. The hub takes it if the claim is owed in full by its book:
its notes are spent, and change beyond $b$ is issued back to her. Then
both sign the fold. If the hub takes the claim and won't fold, her exit
posts the same claim; a claim on the same contract never counts against
itself.

\subsection{The dispute graph}\label{sec:graph}

Every transaction has one input, and every link one output carrying the
value on, plus an anchor for fees; each output has its terminal. The
links and terminals are pre-signed at the opening.

\begin{small}
\begin{verbatim}
contract output
 |- fold                       both, at a withdrawal
 |- default (after T_close)    -> Alice r_A, the hub the rest
 '- request    Alice: one-time signatures on b and on the claim (233 B)
     |- request_timeout (after w_H)     -> Alice: b, both reserves
     '- correct  the hub: one-time signature on b_corr, and 8 slots,
         |       each its signature on correct itself (64 bytes)
         |- correct_timeout (after w_H) -> Alice: b, both reserves
         '- assert   the hub: pi (Lamport-signed digits), pi2, pi3
             |- success (after D2) -> Alice: b_corr's bits, r_A;
             |                         the hub the rest
             '- challenge_assert   Alice: the reveal of pi's labels
                 |- wrongly_challenged_k -> the hub (its proof valid)
                 '- failure (after D1)   -> Alice: b, both reserves
\end{verbatim}
\end{small}

\textbf{Bits.} The terminals pay $b$ in binary: for each bit $i$, an
output of $u\,2^i$ that Alice spends if bit $i$ of her one-time signed
$b$ is set, the hub otherwise. \tx{success} pays $\bcorr$ the same way,
read from the hub's one-time signature in \tx{correct}; a leaf
\tx{equiv\_corr\_i} gives Alice every bit output if she shows two
different signatures under that one-time key (an equivocation).

\textbf{Windows} on Mutinynet (30\,s blocks): $w_H$ = 240 blocks (2
hours) for the hub to correct, and again to assert; $\Delta_2$ = 2,880
(a day) for her challenge; $\Delta_1$ = 240 for the hub's answer.

\textbf{No penalty.} \tx{success} returns her reserve whatever
$\bcorr$. A false claim gains her nothing (she is paid what she is
owed) and costs the hub one dispute; a penalty would also strike an
honest wallet restored from a stale backup.

\section{The exit}\label{sec:exit}

\subsection{The posted claim}\label{sec:claim}

A claim names settled notes of the contract's key, \emph{fallbacks}
(the ids of her spends in flight whose inputs are among them) and $b$.
An exit posts it compact, in exactly 233 bytes:

\begin{center}
\begin{tabularx}{0.9\linewidth}{lL}
\toprule
Bytes & Content \\
\midrule
1 & the counts of notes and fallbacks, 4 bits each \\
46 per note & the id's first 8 bytes, the value (6 bytes), the hub's signature's $R$ \\
32 & the hub's signatures' $s$, half-aggregated into one \\
8 per fallback & the spend id's first 8 bytes \\
the rest & zeros \\
\bottomrule
\end{tabularx}
\end{center}

The contract ($c$, the deposit, the owner) and $b$ (signed apart) are
not in it. Four notes and two fallbacks fit, or three and four.

\textbf{Half-aggregation}~\cite{halfagg1,halfagg2}. For the hub's
signatures $(R_i, s_i)$ on messages $m_i$ under $I$, with weights
$z_0 = 1$ and $z_i = H(I, R_0, m_0, \dots, R_i, m_i)$, the aggregate is
$(R_0, \dots, R_{n-1}, s = \sum z_i s_i)$, valid if
\[
  sG = \sum_i z_i\,(R_i + e_i I),\qquad e_i = H_{\text{BIP340}}(R_i \,\|\, I \,\|\, m_i).
\]
The user computes it from the signatures she holds.

\textbf{Signed in full.} Her one-time (Winternitz) key for the claim,
fixed at the opening, signs all 233 bytes, so Bitcoin's script checks
every byte and no one can swap them. A witness item can't be bound any
cheaper: tapscript has no \tx{OP\_CAT}, so a hash computed in script
can't be compared with signed digits. The cost is 2 digits a byte and 2
witness items a digit: with $b$'s signature and the 2-of-2's, 962
items, inside Bitcoin's 1,000 stack elements. That bounds the claim's
size.

\textbf{Why not shorter tags.} A MAC tag under a key the hub commits to
at the opening would be 8 bytes, but the user can't check it, so a hub
could issue bad tags and later call her notes unissued. A VRF output is
checkable and short, but the payee's tags would come from the hub after
$t$, outside the grant's atomicity. The hub's BIP340 signatures are hers
to check and are released by $t$ with the outputs.

\subsection{What a claim is owed}

The statement counts a posted claim against a set of \emph{cited}
spends, each given with what she signed: inputs, outputs, $T$, the
output pre-signatures, the reveals, and her signature. Let \tx{slots} be
the $R$ values of \tx{correct}'s eight slot signatures.

\begin{enumerate}[nosep]
\item \textbf{Malformed: owed 0.} The bytes aren't a posted claim (count,
  zero padding); a note is named twice; the aggregate signature fails
  under $I$ with the owner as each note's owner; or $b$ exceeds the
  notes' value in units.
\item \textbf{Usable spends.} A cited spend counts if her signature
  verifies on $H(\mathit{sid} \,\|\, C)$, with $C$ recomputed from the
  data given, and its reveal for her key has a final nonce $R$ among
  \tx{slots}.
\item \textbf{Fixed point.} The candidates start as the claim's notes,
  as (short id, value). Repeat until nothing changes: a usable spend
  kills the live candidates among its inputs, and if the claim names it
  a fallback, its output to her (at most one) joins the candidates.
\item \textbf{Owed:} $\min\big(b,\ \lfloor \sum \text{live candidates} / u
  \rfloor\big)$.
\end{enumerate}

\subsection{The hub's statement}

\emph{The claim \tx{request} posts is owed at most the $\bcorr$ that
\tx{correct} signs.} The private input:
\begin{itemize}[nosep]
\item \tx{request}, with its witness, in a block, and the headers from
  that block on, carrying the contract's work $W$ on it;
\item \tx{correct}, with its witness, spending \tx{request}'s output, in
  a block among those headers;
\item the cited spends (at most 8).
\end{itemize}
The statement reads $b$ and the claim from \tx{request}'s signed
digits, and $\bcorr$ and the slots from \tx{correct}'s. It checks
$\mathrm{owed} \le \bcorr$, and outputs the journal. The issuer and the
owner are bound by $D$; every other input is either on chain under the
headers, or signed by her.

\textbf{Faked evidence only hurts the hub.} Evidence can only lower what
she is owed, and evidence failing a check kills nothing. So a hub
feeding the prover bad data leaves her owed more: the bound
$\mathrm{owed} \le \bcorr$ gets harder to prove, not easier. The spends
it can use against her are exactly those she signed, with the $T$ and
reveals she accepted.

\subsection{Reveals: the hub can't use a spend without giving her its
unlock}\label{sec:reveal}

The difficulty: $t$ must reach the user whenever the hub uses her spend
against her. In \tx{correct}'s witness, data no signature covers can be
swapped by whoever gets the transaction mined. A hash lock would need
its hash fixed in the leaf at the opening, before the spend exists. And
script can't check $tG = T$.

\textbf{The construction.} \tx{correct}'s sighash $m_c$ is fixed at the
opening, and \tx{correct} is always posted when the hub disputes her
claim. Each grant carries the hub's adaptor pre-signature
$\sigma' = (R', s')$ under the spend's $T$ on $m_c$; its final nonce is
$R = R' + T$. The user checks
\[
  s'G = R' + eP_H,\qquad e = H_{\text{BIP340}}(R \,\|\, P_H \,\|\, m_c),
\]
before signing the spend, and her signature commits to $\sigma'$ and
$T$. The \tx{correct} leaf has eight slots, each
\[
  \tx{OP\_SIZE 64 OP\_EQUALVERIFY}\ \langle P_H\rangle\ \tx{OP\_CHECKSIGVERIFY}.
\]
To cite a spend, the hub puts the completed signature
$(R, s' + t)$ in a slot; the other slots take any signature of its on
$m_c$.

\textbf{Why she then holds $t$.} Bitcoin checked the slot's signature:
$sG = R + eP_H$. She checked $s'G = R' + eP_H$, with the same $e$,
since $e$ depends on the final $R$. Subtracting, $(s - s')G = T$: she
computes $t = s - s'$ from the chain. The statement needs no sighash and
no $t$: it checks that a slot's $R$ equals the final nonce in the reveal
she signed.

\textbf{What the hub could try.}
\begin{itemize}[nosep]
\item \emph{Sign with another nonce.} It may, but then the statement
  doesn't count the spend.
\item \emph{Use the matching nonce on another message.} The hub knows
  the discrete log of $R$, so it could sign another message with it.
  For the same transaction, the only other messages come from another
  sighash type, whose 65-byte signatures carry their type byte. The
  slot's size check and the statement's both require 64 bytes, the
  default type. The annex, also in the sighash, is excluded because
  \tx{correct} carries her signature from the opening, made without
  one.
\item \emph{One slot for two spends}, choosing $R'_2 = R'_1 + T_1 - T_2$.
  Harmless: her checks force $s'_2 = s'_1 + t_1 - t_2$, and she
  extracts both unlocks from the one signature.
\item \emph{A spend from before the contract opened}, whose grant has no
  reveal on this contract's \tx{correct}. Ruled out: the contract's key
  is fresh, refused by the hub if it has history.
\end{itemize}

\textbf{The user's side.} After \tx{correct}, her wallet looks in the
slots for the final nonce of each in-flight spend's reveal. For each
one found, it computes $t$, completes the outputs, and keeps the payee's
note to hand him. Her own change is already counted in her claim, as a
revived fallback.

\subsection{How much evidence a claim can need}\label{sec:volume}

The user can't make the proof fail. Her lever is to make more evidence
necessary than the statement carries, leaving the hub to correct to a
higher $\bcorr$. Each useful cited spend kills at least one candidate.
A claim has at most 4 notes and 4 fallbacks, and with one output per
owner per spend (the hub's grant policy), each fallback revives at most
one candidate. So at most 8 spends are ever needed: the 8 slots.

\subsection{The hub's side}

On \tx{request}, the hub reads the posted claim and maps each note and
fallback to its book (short id and value). It gathers the evidence: for
each claimed note, the spend that took it, and through fallbacks their
change in turn. It computes what the claim is owed, with the
statement's own function.
\begin{itemize}[nosep]
\item \textbf{Owed in full:} it takes the claim into its book, and
  \tx{request\_timeout} pays her.
\item \textbf{Otherwise:} it posts \tx{correct} with $\bcorr$ = owed and
  the cited spends' completed reveals in the slots, and takes the claim
  at that. Once \tx{request}'s block carries $W$, it checks the
  statement natively, proves it (about 21 minutes on a 4-core server),
  wraps it in Groth16, and asserts.
\end{itemize}
After \tx{success} it sweeps its bit outputs once $\Delta$ has passed.

\section{Costs}

Measured on regtest (the contested exit: a claim of 584 units corrected
to 365, one spend cited, 12 headers):

\begin{center}
\begin{tabularx}{0.9\linewidth}{lL}
\toprule
Item & Cost \\
\midrule
\tx{request} & 11,723 vB (the claim and $b$ one-time signed: 962 witness items) \\
\tx{correct} & 654 vB (8 slots) \\
statement & 1,306k user cycles, 2 segments \\
proof & 973\,s and a 265\,s Groth16 wrap on a 4-core CPU server (the earlier guest, also 2 segments) \\
payment & one round trip, about 40--100\,ms on regtest \\
\bottomrule
\end{tabularx}
\end{center}

\section{Open questions}

\begin{itemize}[nosep]
\item \textbf{The hub and the guest must agree.} Every spend the hub
  records must be one the statement accepts. The hub should check at
  intake with the guest's own functions (the two BIP340 implementations
  are strict, but this is not yet enforced).
\item \textbf{Fees.} Per payment (the issuer's revenue), and liquidity
  per unit time, as in the anchored L2.
\item \textbf{Payees without a contract} hold credit. A payee with one
  is protected through the payer's wallet when the hub withholds $t$:
  the payer extracts it from \tx{correct} and hands the note over.
\item \textbf{Restoring a wallet.} Contract keys derive from the wallet's
  key and each contract's nonce, so a restore needs the nonces (kept in
  the wallet file today).
\item \textbf{The book} is one file rewritten on each change: fine for a
  testnet, not for volume.
\item \textbf{Spends of several owners} carry one reveal per input
  owner's contract; nothing bounds how many such owners a spend has
  beyond its 16 inputs.
\end{itemize}

\begin{thebibliography}{9}
\bibitem{lichen} waxwing, \emph{Lichen}, version 0.5, 2026.
\bibitem{ecash} m0wer and waxwing, \emph{lichen-ecash}, version 0.2, 2026.
\bibitem{halfagg1} K.~Chalkias, F.~Garillot, Y.~Kondi, V.~Nikolaenko,
  \emph{Non-interactive half-aggregation of EdDSA and variants of Schnorr
  signatures}, CT-RSA 2021.
\bibitem{halfagg2} Y.~Chen, Y.~Zhao, \emph{Half-aggregation of Schnorr
  signatures with tight reductions}, ESORICS 2022; J.~Nick, T.~Ruffing,
  draft BIP, \emph{Half-aggregation of BIP 340 signatures}.
\bibitem{bip340} P.~Wuille, J.~Nick, T.~Ruffing, \emph{BIP 340: Schnorr
  signatures for secp256k1}, 2020.
\end{thebibliography}

\end{document}
